% Part: lambda-calculus
% Chapter: syntax
% Section: abbreviated-syntax

\documentclass[../../../include/open-logic-section]{subfiles}

\begin{document}

\olfileid{lam}{syn}{abb}
\olsection{সংক্ষিপ্ত বাক্যগঠন}

\olref[trm]{defn:term}-এ সংজ্ঞায়িত পদগুলি কখনো কখনো লিখতে দুরূহ হয়,
তাই আরও সংক্ষিপ্ত বাক্যগঠন চালু করা সুবিধাজনক। অবশ্যই নিশ্চিত করতে হবে,
সংক্ষিপ্ত সংকেতলিপির পদগুলিও যেন এককভাবে পাঠযোগ্য হয়। বহুল ব্যবহৃত একটি
রূপকে \emph{সংক্ষিপ্ত পদ} বলা হয়; নিয়মগুলি নিচে দেওয়া হলো।

\begin{enumerate}
\item বন্ধনী বাদ দিলে প্রয়োগ বাঁ থেকে ডান দিকে ঘটে। যেমন, $M$, $N$,
  $P$ ও~$Q$ পদ হলে $MNPQ$ হলো $(((MN)P)Q)$-এর সংক্ষিপ্ত রূপ।
\item আবার, বন্ধনী বাদ দিলে ল্যাম্বডা-বিমূর্তনকে সম্ভাব্য সর্বাধিক বিস্তৃত
  পরিসর দেওয়া হয়। যেমন, $\lambd[x][MNP]$-কে
  $(\lambd[x][MNP])$ হিসেবে পড়া হয়।
\item একটি ল্যাম্বডা দিয়ে একাধিক চলরাশির উপর বিমূর্তন করা যায়। যেমন,
  $\lambd[xyz][M]$ হলো $\lambd[x][\lambd[y][\lambd[z][M]]]$-এর সংক্ষিপ্ত
  রূপ।
\end{enumerate}

যেমন,
\[
\lambd[xy][xxyx \lambd[z][xz]]
\]
নিচের পদের সংক্ষিপ্ত রূপ:
\[
(\lambd[x][(\lambd[y][((((xx)y)x)(\lambd[z][(xz)]))])]).
\]

\begin{prob}
সংক্ষিপ্ত পদ $\lambd[g][(\lambd[x][g (x x)])
  \lambd[x][g (x x)]]$-কে পূর্ণ রূপে লেখো।
\end{prob}

\end{document}
